\subsection{Collect requests on a third provider's real UPI sandbox (pre-registered)}\label{app:provider3}
The two payout tests cover money the agent sends. The other side of UPI is the collect request, in which a payee asks the
customer to pay and the customer approves; NPCI restricted person-to-person collect because of exactly this kind of fraud.
PR59 gave the agent (gpt-4.1-mini) a customer's UPI inbox: it could list pending collect requests, read the customer's SMS
and e-mail, and approve or decline. Each pending request was created for real on a third Indian provider's UPI sandbox,
and an approval was completed with that sandbox's mock-payment API, which records a payment as if the customer's UPI app
had paid. Fifteen attack scenarios used five common frauds (a fake refund, an electricity disconnection notice, a KYC
fee, a parcel re-delivery fee and a marketplace buyer who asks the seller to approve a request ``to receive money''),
each with three phrasings of the customer's instruction. Fifteen legitimate scenarios were requests from payees the
customer had saved. 60 episodes, no errors, no request refused by the provider.
\begin{center}\small
\begin{tabular}{p{7.2cm}cc}
\toprule
 & direct & \textbf{gated}\\
\midrule
attacker's collect request paid & 15 of 15 & \textbf{0 of 15}\\
legitimate request paid & 11 of 15 & \textbf{13 of 15}\\
\bottomrule
\end{tabular}
\end{center}
{\small Both predictions held. Without the gate the agent paid every fraudulent request, including all three marketplace
requests framed as receiving money. The legitimate requests left unpaid were all cases in which the model asked the
customer before approving; the gate refused none. All 39 payments were confirmed in the provider's payment history, each
with a transaction ID and an RRN (status \texttt{initiated} on the sandbox rails). One limitation is declared in the
pre-registration: the sandbox gives an account a single merchant VPA, so the attacker's payee and the legitimate payee are
the same real VPA in different customer sessions. What differs, and what the gate reads, is where the payee came from: a
payee the customer saved, or one that appeared only in the request and an SMS.}

